% SGA 3, Expose V, Loop 1 English reconstruction.
% Scope: printed pages 251--252, plus the closing sentence of Lemma 1.1
% at the top of printed page 253; Lemma 1.2 is excluded.
% Authority: Polo--Gille Exp5-13oct24.pdf, component-local pages 1--3
% (combined-reader pages 263--265). Page 3 controls only the first sentence.
% OCR and the Reinhold English Markdown are comparison witnesses only.

\setcounter{footnote}{0}
\chapter*{第 V 篇\\商概形的构造}
\phantomsection\label{exp:sga3-V}
\addcontentsline{toc}{chapter}{第 V 篇. 商概形的构造}
\markboth{第 V 篇. 商概形的构造}%
  {第 V 篇. 商概形的构造}

\begin{center}
  \textit{P. Gabriel 著}
\end{center}

本篇的目的是证明 TDTE~III 中陈述的定理.%
\footnote{\label{note:sga3-V-1}编者注: 即 TDTE~III 的
\SGARefLink{rem:sga3-V-5.1}{定理~5.1}, 5.3,
\SGARefLink{lem:sga3-V-6.1}{6.1}, 6.2 和~\SGARefLink{lem:sga3-V-7.2}{7.2}.
前两个 (分别地, 接下来的两个) 对应于本篇的
\hyperref[thm:sga3-V-4.1]{定理~4.1}
(分别地, \hyperref[thm:sga3-V-7.1]{定理~7.1} 和~\hyperref[thm:sga3-V-8.1]{8.1}).
TDTE~III 的 \SGARefLink{lem:sga3-V-7.2}{定理~7.2} 在第~VI\textsubscript{A}~篇
3.2 和~3.3 中证明.}
若 \(X\) 和 \(T\) 是范畴~\(\mathscr C\) 的两个对象,
则用 \(X(T)\) 代替 \(\operatorname{Hom}_{\mathscr C}(T,X)\).
类似地, 若 \(\varphi:Y\to X\) 是 \(\mathscr C\) 的箭头
(分别地, 若 \(T\) 是其对象), 则 \(\varphi(T)\) 表示
从 \(Y(T)\) 到 \(X(T)\) 的映射 \(g\mapsto\varphi\circ g\):

\SGAThreeDiagram[0.23\textwidth]
  {figures/exp5/p251_phi_of_T_triangle.png}
  {映射 \(\varphi(T):Y(T)\to X(T)\).}
  {fig:sga3-exp5-p251-phi-of-T}
  {Polo--Gille Exp5-13oct24.pdf, 组件第 1 页; 合卷第 263 页; 印刷第 251 页}

同样地, \(T(\varphi)\) 表示从 \(T(X)\) 到 \(T(Y)\) 的映射
\(g\mapsto g\circ\varphi\):

\SGAThreeDiagram[0.23\textwidth]
  {figures/exp5/p251_T_of_phi_triangle.png}
  {映射 \(T(\varphi):T(X)\to T(Y)\).}
  {fig:sga3-exp5-p251-T-of-phi}
  {Polo--Gille Exp5-13oct24.pdf, 组件第 1 页; 合卷第 263 页; 印刷第 251 页}

最后, 若 \(P\) 是概形, 则用 \(\underline P\) 表示~\(P\) 的底集.

作为例外, 本篇不遵循 \SGARefLink{sga3:iv:scholium:4:6:15}{IV~4.6.15}
(\emph{上述文献}, 原版第~227~页顶部) 中关于商的记号约定,
因为这里希望给出一种也适用于并非等价关系的``预等价关系''的商的构造.%
\footnote{\label{note:sga3-V-2}编者注: 即以 \(X\) 为底的群胚;
参见 \hyperref[sec:sga3-V-1]{第~1~节} 末尾的术语.
当 \(\mathscr C\) 为概形范畴时, 以 \(X\) 为底的群胚 \(X_\ast\)
的商 \(p:X\to Y\) 在某些假设下存在
(参见 \hyperref[thm:sga3-V-4.1]{4.1}, \hyperref[lem:sga3-V-6.1]{6.1}
和~\hyperref[thm:sga3-V-7.1]{7.1}). 若此外 \(X_\ast\) 为等价关系,
则在相同假设下, \(p\) 忠实平坦且拟紧, 因而是普遍满态射;
参见\emph{上述文献}.}

\section*{1. \(\mathscr C\)-群胚}
\phantomsection\label{sec:sga3-V-1}
\addcontentsline{toc}{section}{1. \texorpdfstring{\(\mathscr C\)}{C}-群胚}

\phantomsection\label{par:sga3-V-1-a}%
\noindent\textbf{a)} 设 \(\mathscr C\) 是乘积和纤维积都存在的范畴.
先回顾一下, \(\mathscr C\) 中的图

\SGAThreeDiagram[0.39\textwidth]
  {figures/exp5/p252_exact_pair_cokernel.png}
  {平行箭头对及其拟取的余核.}
  {fig:sga3-exp5-p252-exact-pair-cokernel}
  {Polo--Gille Exp5-13oct24.pdf, 组件第 2 页; 合卷第 264 页; 印刷第 252 页}

\noindent 称为\emph{正合的}, 是指 \(p d_0=p d_1\),
且对每个 \(T\in\mathscr C\), 映射 \(T(p)\) 是从 \(T(Y)\)
到 \(T(X_0)\) 中由满足 \(f d_0=f d_1\) 的箭头 \(f:X_0\to T\)
组成的子集的双射. 此时也称 \((Y,p)\) 是 \((d_0,d_1)\) 的\emph{余核},
并记为
\phantomsection\label{eq:sga3-V-1-cokernel}%
\begin{equation*}
  (Y,p)=\operatorname{Coker}(d_0,d_1).
\end{equation*}

\phantomsection\label{par:sga3-V-1-b}%
\noindent\textbf{b)} 例如, 设 \(\mathscr C\) 是赋环空间范畴
\((\mathbf{Esp.An})\). 此时余核 \((Y,p)\) 总是存在, 可描述如下.
将 \(X_0\) 中的点 \(d_0(x)\) 与 \(d_1(x)\) 等同,
并赋予 \(Y\) 商拓扑, 得到 \(Y\) 的底拓扑空间.
于是典范映射 \(\pi:X_0\to Y\) 与 \(d_0,d_1\)
一起诱导 \(Y\) 上环层的一对箭头:

\SGAThreeDiagram[0.66\textwidth]
  {figures/exp5/p252_ringed_space_sheaf_pair.png}
  {直像层之间诱导的箭头对.}
  {fig:sga3-exp5-p252-ringed-space-sheaf-pair}
  {Polo--Gille Exp5-13oct24.pdf, 组件第 2 页; 合卷第 264 页; 印刷第 252 页}

这里 \(\mathcal O_i\) 是 \(X_i\) 的结构层.
选取 \(\pi_*(\mathcal O_0)\) 中由满足 \(\delta_0(s)=\delta_1(s)\)
的截面 \(s\) 组成的子层, 作为 \(Y\) 上的环层.
箭头 \(p\) 以显然的方式定义.

考虑 \((\mathbf{Esp.An})\) 中的下图:%
\footnote{\label{note:sga3-V-3}编者注:
\hyperref[lem:sga3-V-1.1]{引理~1.1} 和~\hyperref[sga3:exp5:lemma:1.2]{1.2}
为新增内容; 在第~\hyperref[sec:sga3-V-5]{5}--\hyperref[sec:sga3-V-9]{9}~节
中多次用到它们.}

\SGAThreeDiagram[0.18\textwidth]
  {figures/exp5/p252_ringed_space_pair_inline.png}
  {箭头对 \((d_0,d_1):X_1\rightrightarrows X_0\).}
  {fig:sga3-exp5-p252-ringed-space-pair}
  {Polo--Gille Exp5-13oct24.pdf, 组件第 2 页; 合卷第 264 页; 印刷第 252 页}

\noindent 设 \((Y,p)\) 为其余核. 若 \(X_0\) 的开子集 \(U\) 满足
\[
  d_0^{-1}(U)=d_1^{-1}(U),
\]
则称其为\emph{饱和的}; 这等价于 \(U=p^{-1}(p(U))\).
此时, 由于 \(Y\) 带有商拓扑, \(p(U)\) 为~\(Y\) 的开子集.

\medskip
\phantomsection\label{lem:sga3-V-1.1}%
\noindent\textbf{引理 1.1.}
\textit{设 \(U\) 是 \(X\) 的饱和开子集, 并令 \(V=p(U)\).
用 \(U_1=d_0^{-1}(U)=d_1^{-1}(U)\) 表示 \(X_1\) 的相应开子集,
用 \(\widetilde d_0,\widetilde d_1\) 和 \(\widetilde p\)
分别表示 \(d_0,d_1\) 在 \(U_1\) 上的限制以及 \(p\) 在 \(U\) 上的限制.
则 \((V,\widetilde p)\) 是 \((\mathbf{Esp.An})\) 中下述箭头对的余核:}%
\footnote{\label{note:sga3-V-4}编者注: 在概形范畴中则不然.
例如, 令 \(S=\operatorname{Spec}(\mathbf C)\),
\(X_0=\mathbf A^2_S=\operatorname{Spec}(\mathbf C[x_1,x_2])\),
令 \(d_1:\mathbf G_{m,S}\times_S\mathbf A^2_S\to\mathbf A^2_S\)
为 \(\mathbf G_{m,S}\) 在 \(\mathbf A^2_S\) 上的位似作用,
\(d_0\) 为到第二个因子的投影, 令 \(U=\mathbf A^2_S\setminus\{m\}\),
其中 \(m=(0,0)\). 则射影直线 \(\mathbf P^1_S\)
在 \((\mathbf{Esp.An})\) 和 \((\mathbf{Sch})\) 中都是
\((\widetilde d_0,\widetilde d_1)\) 的余核,
而 \((d_0,d_1)\) 在 \((\mathbf{Esp.An})\) 中的余核 \(Y\)
为 \(\mathbf P^1_S\) 与点 \(y_0=\{p(m)\}\) 的并.
包含 \(y_0\) 的开子集只有 \(Y\), 且
\(\Gamma(Y,\mathcal O_Y)=\mathbf C\).
若 \(f:\mathbf A^2_S\to T\) 是满足 \(f d_0=f d_1\) 的
\(S\)-概形态射, 而 \(V=\operatorname{Spec}(A)\) 是 \(T\)
中包含 \(t_0=f(y_0)\) 的仿射开子集, 则
\(f^{-1}(V)=\mathbf A^2\), 且环同态 \(A\to\mathbf C[x_1,x_2]\)
经过 \(\mathbf C\) 分解. 这表明
\(S=\operatorname{Spec}(\mathbf C)\) 是范畴 \((\mathbf{Sch}/S)\)
中 \((d_0,d_1)\) 的余核.}

\SGAThreeDiagram[0.40\textwidth]
  {figures/exp5/p252_lemma_1_1_restricted_cokernel.png}
  {\hyperref[lem:sga3-V-1.1]{引理~1.1} 中限制得到的箭头对及其余核.}
  {fig:sga3-exp5-p252-lemma-1-1-restricted-cokernel}
  {Polo--Gille Exp5-13oct24.pdf, 组件第 2 页; 合卷第 264 页; 印刷第 252 页}

验证是容易的.
% Translation statistics: source_lines=162; output_lines=162; inline_math=113; display_math=1; equation=0; equation_star=1; align=0; footnotes=4.
% Review: complete assigned component; pending_TODO=0; all math symbols, labels, references, diagram paths and environments checked; formula prose translated.
